\section{The mean polytope}\label{sec:meanPolytope}

The mean polytope is the set of possible mean parameters given the by a base measure representable distributions.
We define it as
\begin{align*}
    \genmeanset
    = \left\{\contractionof{\probwith,\sencsstatwith,\basemeasurewith}{\selvariable}\wcols\probwith\in\distmeasuredby{\basemeasure} \right\} \, ,
\end{align*}
where we denote by $\distmeasuredby{\basemeasure}$ the set of all probability distributions representable with respect to $\basemeasure$, and by $\sencsstatwith$ the selection encoding of the statistic (see introduction).

%\begin{figure}
%    \begin{center}
%        \input{tikz/mean-polytope/meanset_sketch_generic.tex}
%    \end{center}
%    \caption{Sketch of the statistic mapping the states to the mean parameters.}
%\end{figure}

\subsection{Convex hull}

First of all, we characterize the mean polytope as the convex hull of the slices to the selection encoding.

\begin{lemma}
    For any statistic $\sstat$ and base measure $\basemeasure$, the set of mean parameters is the convex hull of the set
    \begin{align*}
        \genimset
        = \{\sencsstatat{\indexedshortcatvariables,\selvariable}
        \wcols\shortcatindices\in\facstates\ncond\basemeasureat{\indexedshortcatvariables}\neq0 \}
    \end{align*}
    that is $\genmeanset = \convhullof{\genimset}$.
\end{lemma}
\begin{proof}
    This follows from
    \begin{align*}
        \distmeasuredby{\basemeasure}
        = \convhullof{
            \frac{1}{\basemeasureat{\indexedshortcatvariables}} \cdot \onehotmapofat{\shortcatindices}{\shortcatvariables}
            \wcols\shortcatindices\in\facstates\ncond\basemeasureat{\indexedshortcatvariables}\neq0} \, .
    \end{align*}
    and for any vertex of this simplex we have
    \begin{align*}
        \contractionof{\frac{1}{\basemeasureat{\indexedshortcatvariables}} \cdot \onehotmapofat{\shortcatindices}{\shortcatvariables},\sencsstatwith,\basemeasurewith}{\selvariable}
        = \sencsstatat{\indexedshortcatvariables,\selvariable} \, .
    \end{align*}
\end{proof}

%Thus the polytope of mean parameters depends on the base measure only through its support.

\subsection{Faces}

Let us now continue with the investigation of the faces of the mean parameter polytope. % which we define analogously to Def.~2.1 in \cite{ziegler_lectures_2013}. % ! Defined in Ziegler directly with normals

\begin{definition}
    We say that the convex hull of a subset $\genfaceimset\subset\genimset$ is a face of a mean polytope $\genmeanset$, if and only if there is a face normal vector $\canparamofat{\facesymbol}{\selvariable}\in\parspace$ such that
    \begin{align*}
        \genfaceimset = \argmax_{\meanparamwith\in\genimset} \contraction{\meanparamwith,\canparamofat{\facesymbol}{\selvariable}} \, .
    \end{align*}
    We denote the face as $\facesymbol=\convhullof{\genfaceimset}$.
    The set of all faces of the mean polytope $\genmeanset$ is denoted by $\facelatticeof{\genmeanset}$.
\end{definition}

$\facelatticeof{\genmeanset}$ is called a lattice \cite{ziegler_lectures_2013}.


We notice, that each face itself is a convex polytope.
What is more, we can characterize these as mean parameter polytopes with respect to refined base measures to be defined next.

\begin{definition}
    \label{def:faceMeasure}
    The to a face $\facein$ refined base measure is the tensor $\genfacemeasurewith$ with coordinates to $\shortcatindicesin$ by
    \begin{align*}
        \genfacemeasurewith
        = \begin{cases}
              \basemeasureat{\indexedshortcatvariables} & \ifspace \sstatat{\shortcatindices}\in\genfaceimset \\
              0 & \elsetext
        \end{cases} \, .
    \end{align*}
\end{definition}

We now specify the mean parameter polytope to any face using the face measure as a refinement of the base measure.

\begin{lemma}
    \label{lem:faceAsRefinedPolytope}
    For any face $\facein$ we have
    \begin{align*}
        \facesymbol = \meansetof{\sstat,\basemeasureof{\sstat,\facesymbol}} \, .
    \end{align*}
\end{lemma}
\begin{proof}
    For any $\shortcatindicesin$ we have $\genfacemeasureat{\indexedshortcatvariables}\neq0$ if and only if $\sstatat{\shortcatindices}\in\genfaceimset$ and $\basemeasureat{\indexedshortcatvariables}\neq0$.
    Thus we have
    \begin{align*}
        \facesymbol
        &= \convhullof{\genfaceimset} \\
        &= \convhullof{\sencsstatat{\indexedshortcatvariables,\selvariable} \wcols \sstatat{\shortcatindices}\in\genfaceimset\ncond \basemeasureat{\indexedshortcatvariables}\neq 0} \\
        &= \convhullof{\sencsstatat{\indexedshortcatvariables,\selvariable} \wcols \genfacemeasureat{\indexedshortcatvariables}\neq 0} \\
        &= \meansetof{\sstat,\basemeasureof{\sstat,\facesymbol}} \, . \qedhere
    \end{align*}
\end{proof}

Representability of a distribution with respect to face measures is an equivalent condition for the mean parameter of a distribution to be on a face, as we show next.

\begin{lemma}
    \label{lem:faceMeasureRepCondition} % was \label{the:facePolytopeCharacterization} in the report
    If and only if for a distribution $\probwith\in\distmeasuredby{\basemeasure}$ and a face $\facein$ we have
    \begin{align*}
        \contractionof{\sencsstatwith,\probwith,\basemeasurewith}{\selvariable}\in\facesymbol\, ,
    \end{align*}
    then %$\probwith$ is representable with respect to the base measure \red{Did not define representable, reformulate with the support!}
    \begin{align*}
        \probwith \in \distmeasuredby{\genfacemeasure} \, .
    \end{align*}
\end{lemma}
\begin{proof}
    For a distribution $\probwith\in\distmeasuredby{\basemeasure}$ to be in $\distmeasuredby{\genfacemeasure}$, it is enough to show that $\probwith$ is only supported at $\genfacemeasurewith$.
    We have
    \begin{align*}
        \meanparamat{\selvariable}
        = \sum_{\shortcatindices} \probat{\indexedshortcatvariables}\cdot\basemeasureat{\indexedshortcatvariables} \cdot \genstatshortcatencoding \, .
    \end{align*}
    Let now $\canparamofat{\facesymbol}{\selvariable}$ be a face normal to the face $\facesymbol$.
    We then have
    \begin{align*}
        \contraction{\meanparamwith,\canparamofat{\facesymbol}{\selvariable}}
        = \sum_{\shortcatindicesin}
        \probat{\indexedshortcatvariables}\cdot\basemeasureat{\indexedshortcatvariables} \cdot \contraction{\genstatshortcatencoding,\canparamofat{\facesymbol}{\selvariable}} \, .
    \end{align*}

    Now if and only if $\probat{\indexedshortcatvariables}\cdot\basemeasureat{\indexedshortcatvariables}$ is supported only for $\shortcatindices$ with $\sstatat{\shortcatindices}\in\genfaceimset$ we have that
    \begin{align*}
        \contraction{\meanparamwith,\canparamofat{\facesymbol}{\selvariable}}
        = \max_{\meanparamwith\in\genmeanset} \contraction{\meanparamwith,\canparamofat{\facesymbol}{\selvariable}}
    \end{align*}
    which is equal to $\meanparamwith\in\facesymbol$.
    Thus, if and only if $\meanparamwith\in\facesymbol$ then $\probat{\shortcatvariables}$ is only supported at $\shortcatindices$ in the support of $\genfacemeasurewith$.

%    Now, the $\shortcatindices$ with $\genfacemeasureat{\indexedshortcatvariables}=1$ are exactly those, for which the conditions $\facesymbol$ hold straight.
%    If and only if for a $\shortcatindices$ with $\genfacemeasureat{\indexedshortcatvariables}=0$ we have $\probat{\indexedshortcatvariables}>0$, one of the conditions $\facesymbol$ would not hold straight.
%    Thus, if and only if $\probwith$ is representable with respect to $\genfacemeasureat{\shortcatvariables}$, we have $\meanparamat{\selvariable}\in\facesymbol$.
\end{proof}


Let us now investigate tensor network representations of refined base measures, based on the basis encoding $\bencodingof{\sstat}$ of a statistic.
% Vertices
%Vertices of $\genmeanset$ are faces with single elements, that is $\{\meanparamwith\}$.
%By \lemref{lem:faceConvHullPreimage} there must be $\meanparam$ must lie in the image of $\sstatencoding$, since otherwise $\genmeanset$ would be empty.
%The vertex measure is then
%\begin{align*}
%    \basemeasureofat{\sstat,\facesymbol}{\shortcatvariables}
%    = \contractionof{\bencodingofat{\sstat}{\headvariables,\shortcatvariables},\onehotmapofat{\imelementwith}{\headvariables}}{\shortcatvariables}
%\end{align*}
%Here we use that each $\meanparam\in\facesymbol\cap\imageof{\sstatencoding}$ has integer-valued coordinates and denote
%\begin{align*}
%    \onehotmapofat{\imelementwith}{\headvariables} = \bigotimes_{\selindexin} \onehotmapofat{\meanparamat{\indexedselvariable}}{\headvariableof{\selindex}} \, .
%\end{align*}

\begin{theorem}
    \label{the:faceMeasureCharacterization}
    For any face $\facesymbol$ of $\meanset$ we have
    \begin{align*}
        \genfacemeasureat{\shortcatvariables}
        =\contractionof{\sstatccwith,\kcoreofat{\facesymbol}{\headvariables},\basemeasurewith}{\shortcatvariables}
    \end{align*}
    where
    \begin{align*}
        \kcoreofat{\facesymbol}{\headvariables}
        = \sum_{\imelementwith\in\genfaceimset} \onehotmapofat{\imelementwith}{\headvariables} \, .
    \end{align*}
\end{theorem}
\begin{proof}
    For any $\imelementwith\in\genfaceimset$ the tensor
    \begin{align*}
        \hypercoreofat{\imelement}{\shortcatvariables}
        = \contractionof{\sstatccwith,\onehotmapofat{\imelementwith}{\headvariables}}{\shortcatvariables}
    \end{align*}
    is the indicator of the preimage of $\meanparam$ under $\sstatencoding$.
    Since preimages of the elements in $\genfaceimset$ are disjoint, the support of $\hypercoreofat{\imelement}{\shortcatvariables}$ is disjoint and their sum
    \begin{align*}
        \sum_{\imelementwith\in\genfaceimset} \hypercoreofat{\imelement}{\shortcatvariables}
    \end{align*}
    is the indicator of the preimage of $\facesymbol$ under $\sstatencoding$.
    The face measure obeys thus
    \begin{align*}
        \genfacemeasurewith
        &= \contractionof{\left(
                              \sum_{\imelementwith\in\genfaceimset} \hypercoreofat{\imelement}{\shortcatvariables}
        \right), \basemeasurewith}{\shortcatvariables} \\
        &= \sum_{\imelementwith\in\genfaceimset}
        \contractionof{\sstatccwith,\onehotmapofat{\imelementwith}{\headvariables},\basemeasurewith}{\shortcatvariables} \\
        & = \contractionof{\sstatccwith,\kcoreofat{\facesymbol}{\headvariables},\basemeasurewith}{\shortcatvariables} \qedhere
    \end{align*}
\end{proof}

%% Going beyond the report

%%  OLD: restriction to CP
%\begin{definition}
%    The $\cpformat$ rank of a face is
%    \begin{align*}
%        \cprankof{\facesymbol}
%        = \min_{
%            \arbset \wcols \arbset\cap\genimset = \genfaceimset
%            %\{\shortheadindices \wcols \onehotmapof{\shortheadindices}\in\facesymbol\}\subset \arbset \subset \onehotmap(\bigtimes_{\selindexin}[\headdimof{\selindex}]) \ncond
%            %\arbset \cup \{\shortheadindices \wcols \onehotmapof{\shortheadindices}\in\genmeanset/\facesymbol\} = \varnothing
%        }
%        \cprankof{\sum_{v \in \arbset} \onehotmapofat{\imelement}{\headvariables}} \, .
%    \end{align*}
%    By replacing $\cprankof{\cdot}$ with $\bascprankof{\cdot}$ and $\baspluscprankof{\cdot}$ we further define the basis and basis+ CP rank of a face.
%\end{definition}
%
%
%
%The face measures are contraction of the vertex subset encodings with the computation.
%They are \ComputationActivationNetworks{}, when choosing the $\cpformat$ graph with rank at least $\baspluscprankof{\facesymbol}$.
%
%\begin{lemma}
%    For each face of the mean polytope we have
%    \begin{align*}
%        \genfacemeasureat{\shortcatvariables|\varnothing} \in \realizabledistsof{\sstat,\cpformat^{\cprankof{\facesymbol}}} \, ,
%    \end{align*}
%    where $\cpformat^{\cprankof{\facesymbol}}$ is the CP graph with a hidden variable of dimension $\cprankof{\facesymbol}$.
%\end{lemma}
%\begin{proof}
%    We find by definition a set $\arbset$ of basis vectors containing the vertices of the face $\genfacemeasure$ but no further vertices, which has a bas+ $\cpformat$ rank of $\baspluscprankof{\facesymbol}$.
%    We have therefore, that $\normalizationof{\sstatccwith,\onehotmapofat{\arbset}{\headvariables}}{\shortcatvariables}$ is in $\realizabledistsof{\sstat,\cpformat^{\baspluscprankof{\facesymbol}}}$.
%    Further it holds that
%    \begin{align*}
%        \contractionof{\sstatccwith,\onehotmapofat{\arbset}{\headvariables}}{\shortcatvariables}
%        &= \sum_{\shortheadindices\in\arbset \wcols \onehotmapof{\shortheadindices}\in\facesymbol} \contractionof{\sstatccwith,\onehotmapofat{\shortheadindices}{\headvariables}}{\shortcatvariables} \\
%        & \quad \quad + \sum_{\shortheadindices\in\arbset \wcols \onehotmapof{\shortheadindices}\notin\facesymbol} \contractionof{\sstatccwith,\onehotmapofat{\imelement}{\headvariables}}{\shortcatvariables} \\
%        &= \sum_{\shortheadindices\in\arbset \wcols \onehotmapof{\shortheadindices}\in\facesymbol} \contractionof{\sstatccwith,\onehotmapofat{\shortheadindices}{\headvariables}}{\shortcatvariables} \\
%        &= \genfacemeasureat{\shortcatvariables} \, .
%    \end{align*}
%    Here we used, that for $\shortheadindices\in\arbset$ with $\onehotmapof{\shortheadindices}\notin\facesymbol$ is not in the image of $\sstat$ and therefore the contraction of its one-hot encoding with the computation cores vanishes.
%    Thus, $\normalizationof{\sstatccwith,\onehotmapofat{\arbset}{\headvariables}}{\shortcatvariables}$ coincides with the normalized face measure, which is therefore in $\realizabledistsof{\sstat,\cpformat^{\baspluscprankof{\facesymbol}}}$.
%\end{proof}

We now investigate the representation of face measures by \ComputationActivationNetworks{}.

\begin{definition}
    \label{def:faceRepresentability}
    Let $\graph$ be a hypergraph which nodes include $[\seldim]$.
    We say that a face $\facesymbol$ is representable by $\graph$ if and only if there is a set $\arbset$ of basis vectors with $\arbset\cap\genimset = \genfaceimset$ and there is a tensor network $\kcoreof{\graph}$ with respect to the hypergraph $\graph$ such that
    \begin{align*}
        \contractionof{\kcoreof{\graph}}{\headvariables}
        = \sum_{\imelement\in\arbset} \onehotmapofat{\imelement}{\headvariables} \, . % Can allow for arbitrary weights of the non-supported coordinates!
    \end{align*}
    We call any such tensor network a face activating tensor network.
\end{definition}

\begin{lemma}
    \label{lem:faceMeasureRepresentation}
    If any only if a face is representable by a hypergraph $\graph$ we have
    \begin{align*}
        \frac{1}{\contraction{\genfacemeasurewith}} \cdot \genfacemeasurewith
        \in \realizabledistsof{\sstat,\graph,\basemeasure} \, .
    \end{align*}
\end{lemma}
\begin{proof}
    For any $\arbset$ with $\arbset \wcols \arbset\cap\genimset = \genfaceimset$ and tensor network $\kcoreof{\graph}$ respecting the assumptions of \defref{def:faceRepresentability} we have
    \begin{align*}
        \contractionof{\contractionof{\kcoreof{\graph}}{\headvariables},\bencsstatwith}{\shortcatvariables}
        &= \contractionof{\sum_{\imelement\in\arbset}\onehotmapofat{\imelement}{\headvariables},\bencsstatwith}{\shortcatvariables} \\
        &= \contractionof{\sum_{\imelement\in\genfaceimset}\onehotmapofat{\imelement}{\headvariables},\bencsstatwith}{\shortcatvariables} \, .
    \end{align*}
    Here we used in the second equation that only the vertices in $\genimset $ are in the image of $\sstat$.
    It follows, that
    \begin{align*}
        \frac{1}{\contraction{\genfacemeasurewith}} \cdot \genfacemeasurewith
        = \frac{
        \contractionof{\contractionof{\kcoreof{\graph}}{\headvariables},\bencsstatwith}{\shortcatvariables}
        }{
            \contraction{\contractionof{\kcoreof{\graph}}{\headvariables},\bencsstatwith,\basemeasurewith}
        }
    \end{align*}
    and thus $\genfacemeasureat{\shortcatvariables|\varnothing} \in \realizabledistsof{\sstat,\graph,\basemeasure}$.
\end{proof}

Let us now investigate, which normalized face measures can be computed using $\sstat$ and a hypergraph $\graph$.

\begin{example}[Vertices]
    \label{exa:vertexMeasures}
    Vertices $\facesymbol=\{\imelementwith\}$ are proper faces of affine dimension $0$, that is they consist in single vectors $\imelementwith$.
    Since all vertices are in the image $\genimset$, there exists an index tuple $\shortcatindices\in\stateset$ such that $\basemeasureat{\indexedshortcatvariables}=1$ and
    \begin{align*}
        \imelementwith
        = \sencsstatat{\indexedshortcatvariables,\selvariable} \, .
    \end{align*}
    Then $\kcoreofat{\facesymbol}{\headvariables}$ is the one-hot encoding of the by an interpretation map $\indexinterpretation$ assigned index to $\sencsstatat{\indexedshortcatvariables,\selvariable}$, that is
    \begin{align*}
        \kcoreofat{\facesymbol}{\headvariables}
        = \onehotmapofat{\invindexinterpretationat{\sencsstatat{\indexedshortcatvariables,\selvariable}}}{\headvariables} \, .
    \end{align*}
    In particular, the activation core is elementary and the face measure to any vertex is in $\realizabledistsof{\sstat,\elgraph,\basemeasure}$.
\end{example}

While vertices are the minimal non-vanishing faces in the face-lattice (see \cite{ziegler_lectures_2013}), we now show that also the maximal face, namely the polytope itself, is representable with respect to the elementary hypergraph $\elgraph$.

\begin{example}[Maximal face]
    \label{exa:maximalFaceMeasure}
    The maximal face $\facesymbolof{\varnothing}=\genmeanset$ coincides with the mean polytope itself and is given by the choice $\canparamofat{\varnothing}{\selvariable}=\zerosat{\selvariable}$.
    In this case the corresponding activation tensor to the face measure is trivial, that is
    \begin{align*}
        \kcoreofat{\varnothing}{\headvariables}
        = \onesat{\headvariables} \, .
    \end{align*}
    $\kcoreof{\varnothing}$ is elementary and the normalized face measure $\basemeasureof{\sstat,\varnothing}$ to the maximal face is in $\realizabledistsof{\sstat,\elgraph}$.
\end{example}


Extending \exaref{exa:vertexMeasures}, we can provide a coarse estimation of the hypergraph $\graph$ required to decompose $\kcoreof{\facesymbol}$ for generic faces $\facesymbol$.

\begin{lemma}
    \label{lem:CPfaceRepresentationBound} % Sharpened bound
    Any face $\facesymbol$ is representable by a $\cpformat$ graph with hidden rank
    \begin{align*}
        r = \min\left(\cardof{\genfaceimset},\cardof{\genimset}-\cardof{\genfaceimset}+1\right) \, .
    \end{align*}
\end{lemma}
\begin{proof}
    We show the claim by constructing two face activating tensor networks to $\facesymbol$ in a $\cpformat$ hypergraph with hidden rank $\cardof{\genfaceimset}$ and in a $\cpformat$ hypergraph with hidden rank $\cardof{\genimset}-\cardof{\genfaceimset}+1$.
    To show the first representation we enumerate the vertices by a variable $\decvariable$ with dimension $r=\cardof{\genfaceimset}$, i.e. $\genfaceimset = \{v^{\decindex}[\selvariable] \wcols \decindexin\}$.
    Then we define for $\selindexin$ core tensors $\hypercoreofat{\selindex}{\headvariableof{\selindex},\decvariable}$
    \begin{align*}
        \hypercoreofat{\selindex}{\headvariableof{\selindex},\indexeddecvariable} = \onehotmapofat{\imelementofat{\decindex}{\indexedselvariable}}{\headvariableof{\selindex}} \, .
    \end{align*}
    Then we have
    \begin{align*}
        \contractionof{\{\hypercoreofat{\selindex}{\headvariableof{\selindex},\decvariable}\wcols\selindexin\}}{\headvariables}
        =\sum_{\imelement\in\genfaceimset} \onehotmapofat{\imelement}{\headvariables}
    \end{align*}
    and thus have found an activation tensor network in a $\cpformat$ graph with hidden rank $\cardof{\genfaceimset}$ representing the face $\facesymbol$.

    We continue with the second representation, for which we enumerate the set $\genimset/\genfaceimset$ by $v^{\decindex}[\selvariable]$ where $\decindex\in[\cardof{\genimset}-\cardof{\genfaceimset}]$.
    We define variable $\decvariable$ with dimension $r=\cardof{\genimset}-\cardof{\genfaceimset}+1$ and define for $\selindexin$ core tensors
    \begin{align*}
        \hypercoreofat{\selindex}{\headvariableof{\selindex},\indexeddecvariable}
        =
        \begin{cases}
            -\onehotmapofat{\imelementofat{\decindex}{\indexedselvariable}}{\headvariableof{\selindex}} & \ifspace \decindex < \cardof{\genimset}-\cardof{\genfaceimset} \\
            \onesat{\headvariableof{\selindex}} & \ifspace \decindex = \cardof{\genimset}-\cardof{\genfaceimset}
        \end{cases} \, .
    \end{align*}
    We then have
    \begin{align*}
        \contractionof{\{\hypercoreofat{\selindex}{\headvariableof{\selindex},\decvariable}\wcols\selindexin\}}{\headvariables}
        &=\onesat{\headvariables} - \sum_{\imelement\in\genimset/\genfaceimset} \onehotmapofat{\imelement}{\headvariables} \\
        &= \sum_{\imelement\in\genfaceimset} \onehotmapofat{\imelement}{\headvariables}
        + \sum_{\imelement\in\left(\bigtimes_{\selindexin}[\seldimof{\selindex}]\right)/\genimset} \onehotmapofat{\imelement}{\headvariables} \, .
    \end{align*}
    We have thus found a face activating tensor network for $\facesymbol$ in a $\cpformat$ format with hidden rank $r=\cardof{\genimset}-\cardof{\genfaceimset}+1$.
\end{proof}

\subsection{Examples of polytopes with elementary representable faces}

We first show in the following example, that all faces of a hypercube are representable by elementary activation tensors.

\input{examples/mean-polytope/hypercube_faces}

\begin{example}[Simplices]
    \label{exa:simplexFaces}
    The $(\seldim-1)$-dimensional standard simplex $\meansetof{\triangle,\seldim-1}$ in $[0,1]^{\seldim}$ is the convex hull of the vertex sets
    \begin{align*}
        \imset = \{\onehotmapofat{\selindex}{\selvariable} \wcols\selindexin\} \, . % selvariable
    \end{align*}
    %where $\selvariable$ takes values in $[\seldim-1]$.
    The simplex is the mean polytope of the $\seldim$-dimensional statistic with
    \begin{align*}
        \sstatcoordinateof{\selindex} \defcols [\seldim] \rightarrow [2] \quad, \quad
        \sstatcoordinateofat{\selindex}{\secselindex}
        = \begin{cases}
              1 & \ifspace \secselindex = \selindex \\
              0 & \ifspace \secselindex \neq \selindex
        \end{cases} \, .
    \end{align*}
    This is sometimes referred to as the probability simplex, since it represents all distributions of a discrete variable with $\seldim$ states.

    The face lattice of the simplex consists is enumerated by subsets of $[\seldim]$ as
    \begin{align*}
        \facelatticeof{\meansetof{\triangle,\seldim-1}}
        = \{\facesymbolof{\arbset}\wcols\arbset\subset[\seldim]\}
    \end{align*}
    where the faces are
    \begin{align*}
        \facesymbolof{\arbset}
        = \convhullof{\onehotmapofat{\selindex}{\selvariable} \wcols\selindex\in\arbset}
    \end{align*}
    and therefore itself $(\cardof{\arbset}-1)$-dimensional standard simplices.
    The partial order of the faces coincides with the inclusion order of subsets $\arbset$.

%    Furthermore, each face $\facesymbolof{\arbset}$ is cube-like since
%    \begin{align*}
%        \facesymbolof{\arbset}
%    = \meansetof{\triangle,\seldim-1} \cap \facesymbolof{[\seldim]/\arbset,0_{[\seldim]/\arbset}}
%    \end{align*}
%    where $\facesymbolof{[\seldim]/\arbset,0_{[\seldim]/\arbset}}$ are faces of the hypercube following the notation of \exaref{exa:hypercubeFaces}.
\end{example}

\subsection{Partition into Relative Interiors of Faces}

Let us now introduce relative interiors, which enables us to find disjoint partitions of the mean polytope.

\begin{definition}[Relative Interior]
    \label{def:relativeInterior}
    The relative interior $\sbinteriorof{\arbset}$ of an arbitrary set $\arbset\subset\parspace$ is the interior of $\arbset$ in the affine subspace $\convhullof{\arbset}$.
\end{definition}

\begin{lemma}
    \label{lem:relativeInteriorPolytopePartition}
    Any polytope is a disjoint union of the relative interiors of its faces, that is
    \begin{align*}
        \genmeanset = \bigcup_{\facein}^{\cdot} \sbinteriorof{\facesymbol} \, .
    \end{align*}
\end{lemma}
\begin{proof}
    For any $\meanparam\in\genmeanset$ we find a face such that $\meanparam\in\facesymbol$.
    If $\meanparam\notin\sbinteriorof{\facesymbol}$, then there is a face ${\tilde{\facesymbol}}\subset\facesymbol$ of smaller affine dimension such that $\meanparam\in\tilde{\facesymbol}$.
    We then redo the above argument for $\facesymbol=\tilde{\facesymbol}$ until we arrive at a face with $\meanparam\in\sbinteriorof{\facesymbol}$.
    We always reach such a face such a face, since the affine dimension of the faces is decreasing and lower bounded by $0$, and the faces of affine dimension $0$ are vertices for which $\facesymbol=\sbinteriorof{\facesymbol}$.
\end{proof}

\begin{definition}
    To each $\meanparam\in\genmeanset$ we denote the unique face $\facesymbol$ with $\meanparam\in\sbinteriorof{\facesymbol}$ by $\faceto{\meanparam}$.
\end{definition}